% Generato da scripts/build_thesis_tables.py dai CSV versionati: non modificare a mano.
% Richiede \usepackage{booktabs}.
\begin{table}[htbp]
\centering
\small
\caption{Tasso di rovina col floor sul calendario sintetico, per $B_0/F$ e orizzonte}
\label{tab:floor_sintetico}
\begin{tabular}{rrr}
\toprule
$B_0/F$ & 150 date & 1000 date \\
\midrule
10 & 0,5910 & 0,8914 \\
31,6 & 0,0501 & 0,5963 \\
100 & 0,0000 & 0,0331 \\
220 & 0,0000 & 0,0000 \\
316,2 & 0,0000 & 0,0000 \\
1000 & 0,0000 & 0,0000 \\
3162,3 & 0,0000 & 0,0000 \\
10\,000 & 0,0000 & 0,0000 \\
22\,000 & 0,0000 & 0,0000 \\
31\,622,8 & 0,0000 & 0,0000 \\
100\,000 & 0,0000 & 0,0000 \\
316\,227,8 & 0,0000 & 0,0000 \\
1\,000\,000 & 0,0000 & 0,0000 \\
3\,162\,277,7 & 0,0000 & 0,0000 \\
10\,000\,000 & 0,0000 & 0,0000 \\
\bottomrule
\end{tabular}

\smallskip
\begin{minipage}{\linewidth}
\footnotesize
Calendario a una partita per data; quota 2,5, EV vero $-$0,02, edge stimato 0,03, $\lambda$ = 0,25; 10\,000 repliche; floor attivo.\par
Esperimento C, configurazione (a), $B_0/F = 50$: tempo di rovina medio 2480,8 date, mediano 1427,0; fattore di accelerazione 13,3 (nota 2.9: 23).\par
Fonte: \texttt{results/us\_c6\_3\_floor\_synthetic.csv} (SHA256 \texttt{0995031b6f8f}).
\end{minipage}
\end{table}
